\subsection*{All-prime arithmetic normalizer recognition on the actual genus-two target}

\paragraph{Theorem.}
Let $X_5$ be the smooth proper curve $y^2=1-x^5$, let $O$ be its
unique point at infinity, and let $U_5=X_5\setminus\{O\}$.
Over every number field $K$, and for every rational prime $\ell$,
every nontrivial closed procyclic subgroup
\[
 A\subset P=\pi_1(U_{5,\overline K})^{(\ell)}
\]
whose arithmetic normalizer has open image in $G_K$ is contained
in actual cusp inertia. The result concerns this specified
genus-two curve; arbitrary-cover and full profinite A.22 remain
unresolved.

\paragraph{1. Reduction to the four split CM lines.}
F.3.21 proves the assertion for $\ell=5$ or
$\ell\not\equiv1\pmod5$. We now take $\ell\equiv1\pmod5$.
In particular $\ell\ne2,5$. Write $E=\mathbf Q(\zeta_5)$,
$B=J_{X_5}$, and $V=V_\ell B$. F.3.19 proves absolute
simplicity, and the dimension argument of F.3.21 gives
\[
 \operatorname{End}^0_{\overline{\mathbf Q}}B=E.
\]
All these endomorphisms are defined over $E$. After extending
$K$ finitely, we may suppose $E\subset K$; the normalizer's
image stays open, and the geometric group $P$ is unchanged.
Fix an embedding $\iota:E\hookrightarrow\mathbf Q_\ell$.
The automorphism
\[
 t:(x,y)\longmapsto(\zeta_5x,y)
\]
splits $V$ into four lines $V_a$, $a=1,2,3,4$, on which $t$
has eigenvalue $\iota(\zeta_5)^a$. Their $G_E$-characters are
denoted $\chi_a$.

By Faltings, the endomorphism algebra on every open arithmetic
subgroup is $E\otimes\mathbf Q_\ell=\mathbf Q_\ell^4$.
Semisimplicity then implies that its one-dimensional invariant
subspaces are precisely these four lines. Indeed a stable
line has an equivariant projector of rank one; the only
rank-one idempotents of the displayed rank-one module are
the four primitive coordinate projectors. This also excludes
coincidence of two characters on an open subgroup.

Let $u$ generate $A$. The actual finite-cover detection and
weight comparison of F.3.21 apply before its nonsplit-prime
exclusion. Thus either the first nonzero Magnus term has
degree one, or the detected cover-homology line lies in its
boundary kernel. In the latter case the exact boundary
argument and full-normalizer character transfer of F.3.21
give $\theta=\chi_\ell$ on an open subgroup and prove the
theorem. It remains to rule out degree one. Its nonzero
linear term lies on one of the four $V_a$.

\paragraph{2. All character relations, proved through local Hodge--Tate maps.}
The principal polarization is preserved by $t$. The Weil
pairing therefore pairs $V_a$ only with $V_{5-a}$, and gives
the exact character equalities
\[
 \chi_1\chi_4=\chi_2\chi_3=\chi_\ell.
 \tag{X5.1}
\]
There are no further finite-image monomial relations on any
open subgroup of $G_E$.

Here is a direct proof of the latter assertion, including
the local realization used. At every prime above $\ell$,
$X_5$ has good reduction. One smooth chart is
$y^2=1-x^5$; its chart at infinity is
$v^2=s^6-s$, where $s=1/x$ and $v=y/x^3$.
Their derivatives and $\ell\ne2,5$ give smoothness; they
glue to the proper double cover of $\mathbf P^1$ branched
at the five simple roots and infinity. Its Jacobian is an
abelian scheme. Tate's functorial Hodge--Tate decomposition
for its $\ell$-divisible group, as checked in B.1.4 from
Tate~1967, \S4.1, Theorem~3 and Corollary~2, identifies the
exponent-one summand of $V\otimes C_\ell$ with
$\operatorname{Lie}B\otimes C_\ell(1)$.
The forms $dx/y,x\,dx/y$ show that $t$ on the tangent
space has eigenvalues $\zeta_5,\zeta_5^2$.

The four embeddings of $E$ into $\mathbf Q_\ell$ give
exponent-one index sets
\[
 \{1,2\},\quad \{2,4\},\quad \{4,3\},\quad \{3,1\}.
 \tag{X5.2}
\]
Every finite extension of $E$ has a place over each of these
four places; restriction retains these exponents. A
finite-image character has exponent zero after its finite
trivializing extension, hence already has exponent zero.
For a finite-image product $\prod_a\chi_a^{n_a}$ the four
rows of (X5.2) give
\[
 n_1+n_2=n_2+n_4=n_4+n_3=n_3+n_1=0.
\]
Solving gives
$(n_1,n_2,n_3,n_4)=(b,-b,-b,b)$. By (X5.1) that
product is identically one. This proves the assertion on
every open subgroup, without a CM-image openness theorem.

Choose a symplectic basis
\[
 x_1\in V_1,\quad y_1\in V_4,\quad
 x_2\in V_2,\quad y_2\in V_3,
 \qquad \omega=[x_1,y_1]+[x_2,y_2].
\]
The characters of the diagonal action are
$(\alpha_1,\chi_\ell/\alpha_1,
  \alpha_2,\chi_\ell/\alpha_2)$, with
$\alpha_1=\chi_1,\alpha_2=\chi_2$.
It follows from the proved relation lattice that the
$\chi_\ell^3$-weight part of the degree-six free Lie algebra
is exactly
\[
 \bigoplus_{j=0}^3 L_6(j,j,3-j,3-j).
 \tag{X5.3}
\]
Here the entries are the multiplicities of
$x_1,y_1,x_2,y_2$. The assertion remains true on every open
arithmetic subgroup. Every other character appearing in the
finite-dimensional degree-six derivation space has quotient
by $\chi_\ell^3$ of infinite image.

\paragraph{3. Four actual conjugate subgroups and one common arithmetic group.}
Because $U_5$ is defined over $\mathbf Q$, work in the
actual arithmetic quotient
\[
 1\longrightarrow P\longrightarrow\Pi_{\mathbf Q}
   \xrightarrow{p}G_{\mathbf Q}\longrightarrow1.
\]
The group over $K$ is $p^{-1}(G_K)$, after compatible
geometric base points. Let $N$ be the normalizer of $A$
there, and $H=p(N)$, which is open in $G_{\mathbf Q}$.
The first Magnus line of $A$ is some $V_a$.

For $b\in(\mathbf Z/5)^\times$, choose
$\gamma_b\in G_{\mathbf Q}$ with
$\gamma_b(\zeta_5)=\zeta_5^b$, and choose an actual lift
$\widetilde\gamma_b\in\Pi_{\mathbf Q}$, taking
$\gamma_1=\widetilde\gamma_1=1$. Set
\[
 A_b=\widetilde\gamma_b A\widetilde\gamma_b^{-1},\qquad
 N_b=\widetilde\gamma_b N\widetilde\gamma_b^{-1},\qquad
 H_b=\gamma_bH\gamma_b^{-1}.
 \tag{X5.4}
\]
Each $N_b$ normalizes $A_b$ and maps onto the open subgroup
$H_b$. The formula
$\gamma_b t\gamma_b^{-1}=t^b$ gives
\[
 \gamma_b(V_a)=V_{ab^{-1}}.
\]
The four first Magnus lines of $A_b$ are therefore exactly
the four distinct CM lines. Changing a chosen arithmetic
lift by an element of $P$ changes only geometric conjugacy,
which does not change its first linear term.

Choose a finite Galois extension $L/\mathbf Q$ such that
\[
 G_L\subset G_E\cap\bigcap_b H_b.
 \tag{X5.5}
\]
Such an $L$ exists by taking the normal core of that open
intersection. In particular $G_L$ is contained in each
$G_{\gamma_b(K)}$, since $H_b$ is; all four arithmetic
covering groups can now be restricted to the same field.
For each $b$, the closed subgroup
$N_b\cap p^{-1}(G_L)$ maps \emph{onto} $G_L$.
No simultaneous section is chosen or needed.

\paragraph{4. A special leading representative of the same outer action.}
Let $G_{\mathbf Q}(6)$ be the inverse image of the sixth
weight-filtration subgroup under the outer action for
$U_5$. For any
$\sigma\in G_L\cap G_{\mathbf Q}(6)$, take an automorphism
representative trivial through degree six in the rational
Malcev Lie algebra of $P$. Its action on the boundary loop
$z$ is conjugation, since its action on $V$ and hence on
$\chi_\ell$ is trivial.

One may make this representative fix $z$ while retaining
that filtration. Here are the details needed in genus two.
Write $\log z=\omega+$ higher terms and let $W$ be the
logarithm of a conjugator between its two images. Comparison
in degrees three and four first gives
$[W_1,\omega]=[W_2,\omega]=0$. The centralizer
of $\omega$ in the homogeneous free Lie algebra is
$\mathbf Q_\ell\omega$: a nonzero homogeneous element has
a least Lyndon word; commutation with the least word
$x_1y_1$ of $\omega$ forces these two words to commute.
Two commuting words are powers of a common word, whereas
a Lyndon word is not a proper power. Thus the degree must
be two; subtracting its multiple of $\omega$ then rules
out any remaining leading word. Consequently $W_1=0$,
and the only possible lower term is $W_2=c\omega$.
Multiplying the conjugator on the right by
$\exp(-c\log z)$ leaves its conjugation of $z$ unchanged
and removes that term. For the adjusted conjugator,
comparison in degrees five, six and seven successively
gives $W_3=W_4=W_5=0$. The required adjustment starts in degree
six.

The resulting degree-six derivation $D_\sigma$ is special:
$D_\sigma(\omega)=0$. It is uniquely determined by the
outer leading class, since two such representatives differ
by $\operatorname{ad}w$ with $w\in L_6$ and
$[w,\omega]=0$, forcing $w=0$. The construction is
$\operatorname{GSp}(V)$-equivariant by this uniqueness.
It is only a rational leading representative; the Galois
element $\sigma$ itself remains an actual element of $G_L$.

For each of the four subgroups in (X5.4), choose a lift of
this \emph{same} $\sigma$ in its restricted $N_b$. Its
conjugation exponent on a generator $u_b$ is one, because
the first linear term of $u_b$ is nonzero and $\sigma$ acts
trivially on $V$. Thus this lift centralizes $u_b$.
The boundary-fixing representative of $\sigma$ differs
from that lift by an inner automorphism. If
$\log u_b=a_bv_b+$ higher terms, the conjugacy equation is
\[
 \widetilde\sigma(\log u_b)
  =\exp(\operatorname{ad}W_b)(\log u_b).
\]
Its degree-two part gives $(W_b)_1=c_bv_b$. Right
multiplication of the conjugator by
$\exp(-(c_b/a_b)\log u_b)$ removes this first term without
changing its conjugation action. The free-Lie centralizer
of the generator $v_b$ then successively gives
$(W_b)_2=\cdots=(W_b)_5=0$. In degree seven one obtains
\[
 D_\sigma(v_b)\in[v_b,L_6].
 \tag{X5.6}
\]
The four separate lifts of the one $\sigma$ have therefore
given four simultaneous constraints on its unique special
derivation $D_\sigma$. Nonprimitive linear coefficients
$a_b$ cause no problem, since this conjugator adjustment
is made over $\mathbf Q_\ell$.

\paragraph{5. Projection inside the actual arithmetic image.}
Nakamura~1996, Corollary~4.2 and table~4.3, p.~115,
Lemma~4.5, p.~116, and (4.7.1)--(4.7.2), p.~117, supply
the nonzero arithmetic degree-six quotient
\[
 q_6:\operatorname{gr}_6\Gamma^{\mathrm{univ}}_{2,1}
       \longrightarrow \mathcal G_{2,1}^{,6},
\]
on which the similitude acts by $\chi_\ell^3$.
The actual arithmetic image for $U_5/L$ maps onto this
quotient rationally. The source expressly retains this
graded image under finite base extension, so the field
chosen in (X5.5) does not remove the nonzero quotient.

Let $M$ be the image of $G_L\cap G_{\mathbf Q}(6)$ in the
degree-six graded space, embedded into special derivations
by Step~4. It is a compact additive group in a finite
dimensional $\mathbf Q_\ell$-space, hence a closed
$\mathbf Z_\ell$-module contained in a finite lattice.
It is invariant under conjugation by $G_L$ and every
member satisfies all four equations (X5.6).
Choose $m\in M$ with $q_6(m)\ne0$.

The action on this derivation space is diagonal in the
four CM characters. By Step~2 there exists
$\tau\in G_L$ for which $\chi_\ell(\tau)^3$ differs from
every other character value occurring in that finite
space. To justify existence, the kernel of any of the
finitely many nontrivial quotient characters has empty
interior: an open kernel would give finite image, which
Step~2 excludes. These kernels are closed, and the Baire
property of the compact group $G_L$ excludes their finite
union.

For $T=\operatorname{Ad}(\tau)$ use the polynomial
spectral projector $e(T)$ onto eigenvalue
$\chi_\ell(\tau)^3$. After multiplying by a power of
$\ell$, its coefficients are in $\mathbf Z_\ell$; hence
\[
 m_0=\ell^r e(T)m\in M,\qquad
 q_6(m_0)=\ell^r q_6(m)\ne0.
 \tag{X5.7}
\]
The defining image property of $M$ lifts $m_0$ to an actual
$\sigma_0\in G_L\cap G_{\mathbf Q}(6)$. Its leading
derivation is special, has character $\chi_\ell^3$, and
still satisfies all four tangency constraints. The same
$q_6$ is used throughout; it is the arithmetic projection
of the moduli problem for this curve, independent of the
four chosen normalizer lifts.

\paragraph{6. The complete finite Lie-algebra calculation.}
Put $a=x_1,b=y_1,c=x_2,d=y_2$ and order the alphabet
$a<b<c<d$. Let $b_w$ be the standard Lyndon bracket.
Tangency means
\[
 D(v)=[A_v,v]\quad(v=a,b,c,d),\qquad A_v\in L_6.
\]
The map $A_v\mapsto[A_v,v]$ is injective in degree six.
Equivariance and Step~2 therefore put each $A_v$ in the
balanced space (X5.3). Its four summands have respective
dimensions $3,30,30,3$, as given by their full Lyndon bases.

The equation $D([a,b]+[c,d])=0$ splits into five
multidegrees $(j,j,4-j,4-j)$, $j=0,\ldots,4$. For each
column labelled $(v,w)$ its polynomial is, respectively,
\[
 [[b_w,a],b],\quad [a,[b_w,b]],\quad
 [[b_w,c],d],\quad [c,[b_w,d]].
\]
The row indexed by a degree-eight Lyndon word is its
associative-word coefficient. Restriction to these rows
is injective on the homogeneous free Lie algebra: in a
nonzero combination of Lyndon brackets its least Lyndon
word has the unchanged nonzero coefficient. Thus these
matrices impose the full specialness equation.

The complete certificate
\texttt{check\_genus\_two\_tangent\_space.py} constructs
every column and every row from these definitions. It
gives
\[
\begin{array}{c|rrrrr}
 j&0&1&2&3&4\\\hline
 \text{rows}&8&140&312&140&8\\
 \text{columns}&6&66&120&66&6\\
 \text{rank over }\mathbf Q&6&66&119&66&6\\
 \text{nonzero minor modulo }101&100&1&99&100&100.
\end{array}
\tag{X5.8}
\]
The output includes every selected row and column index.
The nonzero minors prove the rational lower ranks, hence
the same ranks over every $\mathbf Q_\ell$. In the middle
block the following exact nonzero kernel supplies the
matching upper rank:
\[
 A_a=A_b=b_{abcdcd}+2b_{acdbcd}+b_{acdcdb},\qquad
 A_c=A_d=b_{ababcd}+b_{abacdb}.
 \tag{X5.9}
\]
The script checks this equality in the full associative
word expansion, not only in the selected rows.

Equivalently, putting
\[
 z_1=[a,b],\quad z_2=[c,d],\quad C=[z_1,z_2],
 \quad A_1=[z_2,C],\quad A_2=-[z_1,C],
\]
the entire space is the line spanned by
\[
 D_*(a)=[A_1,a],\quad D_*(b)=[A_1,b],\quad
 D_*(c)=[A_2,c],\quad D_*(d)=[A_2,d].
 \tag{X5.10}
\]
The coefficients in (X5.9) give $-D_*$. The identity
$C=[z_1,\omega]$ proves that every value of $D_*$ lies
in the Lie ideal $(\omega)$.

\paragraph{7. The actual universal quotient kills this entire line.}
For $r=0,1$, put
$\Pi_r^{\mathrm{mod}}=\pi_1(M_{2,r})$, write $p_r$ for
its arithmetic projection and $\Phi_r$ for universal outer
monodromy, and set
\[
 F_r(m)=\Phi_r^{-1}(\Gamma_{2,r}(m)),\qquad
 Q_r(m)=\overline{\mathbf Q}^{\,p_r(F_r(m))}.
\]
Takao, RIMS~1709 (November~2010), Theorem~3.5, p.~21,
proves $Q_0(m)=Q_1(m)$ for every $m$. This theorem includes
the last-puncture step in genus two; the older NTU~1994
affine-only assertion is not substituted for it.

The actual morphism $M_{2,1}\to M_{2,0}$ and the quotient
of the surface group obtained by filling the puncture
give the commutative diagram
\[
\begin{array}{ccc}
 \operatorname{gr}_6\Gamma^{\mathrm{univ}}_{2,1}
 &\longrightarrow&
 \operatorname{gr}_6\Gamma^{\mathrm{univ}}_{2,0}\\
 q_{6,1}\downarrow&&\downarrow q_{6,0}\\
 \operatorname{Gal}(Q_1(7)/Q_1(6))\otimes\mathbf Q_\ell
 &\xrightarrow{\ =\ }&
 \operatorname{Gal}(Q_0(7)/Q_0(6))\otimes\mathbf Q_\ell.
\end{array}
\tag{X5.11}
\]
The lower arrow is induced by equal subfields of
$\overline{\mathbf Q}$. The two arithmetic projections
agree on each element of $F_1(6)$, which proves
commutativity. The maps descend to the graded monodromy
images because an element in $F_r(7)$ acts trivially on
$Q_r(7)$. Clearing denominators, and then any torsion in
an integral graded quotient, proves the same statement
for rational classes. On the affine side $q_{6,1}$ is
the exact map (4.7.1) of Nakamura~1996.

Here the precise Lie map needed for filling is as follows.
The inclusion $U_5\hookrightarrow X_5$ induces the quotient
of $P$ by the closed normal subgroup generated by its
boundary loop $z$. Its associated graded Lie algebra is
generated by degree one, and $z$ has initial term
$\omega$. The induced graded map therefore factors as
the canonical surjections
\[
 L(V)\longrightarrow L(V)/(\omega)
       \longrightarrow
 \operatorname{gr}\pi_1(X_{5,\overline{\mathbf Q}})^{(\ell)}
       \otimes\mathbf Q_\ell.
\]
This factorization, proved directly from the actual
boundary relator, is sufficient; no injectivity of the
second arrow is used. Since every $D_*(v)$ belongs to
$(\omega)$, all its values vanish in the proper graded
Lie algebra. Thus the top horizontal arrow kills
\[
 \mathbf Q_\ell[D_*]\cap
     \operatorname{gr}_6\Gamma^{\mathrm{univ}}_{2,1}.
\]
Diagram (X5.11) proves that $q_6$ kills this intersection.
Step~6 places the \emph{actual} $m_0$ of (X5.7) in this
intersection, contradicting $q_6(m_0)\ne0$.
This excludes the remaining degree-one case.

\paragraph{8. Return to the actual subgroup.}
With degree one excluded, the detecting-cover line lies
in the boundary kernel by the weight argument of F.3.21.
After a finite extension splitting those boundary points,
the full normalizer acts on $A$ by exactly $\chi_\ell$.
Nakamura~1994, Theorem~2.1.1, p.~94, then puts $A$ in
actual cusp inertia. Every operation involving formal
conjugators above was used only to derive a contradiction
for an actual arithmetic element; the final subgroup and
recognition take place in the original pro-$\ell$ group.

\paragraph{An exact consequence for the exhibited cover.}
The actual morphism $q:Y_{20}\to U_5$ constructed in
F.3.19 induces compatible arithmetic fundamental-group
maps. For a procyclic subgroup of the geometric group of
$Y_{20}$ whose image under $q_*$ is nontrivial, its
normalizer maps to a normalizing subgroup with the same
open arithmetic image. The theorem places that image
in cusp inertia of $U_5$. Its image in $H_1(U_5)$ is
therefore zero, since $U_5$ has one puncture. Consequently
no subgroup with open arithmetic normalizer can have
nonzero image under the actual map
$H_1(Y_{20})\to V_\ell J_{X_5}$ induced by $q$.
This closes that particular CM-surface detecting target.
It gives no elliptic or genus-two detecting morphism for
an arbitrary finite cover in the full A.22 assertion.
